% Generated draft: requires mathematical and source review.
\documentclass[11pt]{article}
\usepackage[margin=25mm]{geometry}
\usepackage{fontspec}
\setmainfont{Latin Modern Roman}
\usepackage{amsmath,amssymb,mathtools,unicode-math}
\setmathfont{Latin Modern Math}
\usepackage{enumitem,needspace,xurl,hyperref}
\setlength{\parindent}{0pt}
\setlength{\parskip}{5pt}
\setlength{\emergencystretch}{3em}
\newcommand{\N}{\mathbb{N}}
\newcommand{\Z}{\mathbb{Z}}
\newcommand{\Q}{\mathbb{Q}}
\newcommand{\R}{\mathbb{R}}
\newcommand{\C}{\mathbb{C}}
\newcommand{\F}{\mathbb{F}}
\newcommand{\A}{\mathbb{A}}
\newcommand{\PP}{\mathbb{P}}
\newcommand{\E}{\mathbb{E}}
\newcommand{\eps}{\varepsilon}
\newcommand{\dd}{\,\mathrm{d}}
\begin{document}
\section*{UnsolvedMath ID: 30002899}
\textbf{Source problem code:} catalogue entry
\subsection{Definitions and mathematical statement}
For a group \(G\) and \(K\in\{\R,\C\}\), the group ring \(K[G]\) consists of finite sums \(a=\sum_{g\in G}a_g g\), with multiplication induced by the group operation. Its standard involution is
\[
 a^*=\sum_{g\in G}\overline{a_g}\,g^{-1},
\]
where conjugation is the identity for \(K=\R\). An element is Hermitian if \(a=a^*\), and is a finite sum of Hermitian squares if
\[
 a=\sum_{j=1}^{m}b_j^*b_j,
 \qquad m\in\Z_{\geq 0},\quad b_j\in K[G].
\]
For a nonzero commutative coefficient ring \(S\), the cohomological dimension \(\operatorname{cd}_S(G)\) is the projective dimension of the trivial \(S[G]\)-module \(S\). Thus \(\operatorname{cd}_S(G)\leq2\) means that this module admits a projective resolution of length at most two.

The supplied question is: \emph{Does Helton's sum-of-squares theorem hold for group rings of groups of cohomological dimension at most two?}

A possible precise interpretation, presented only to identify the missing choices, would be the following assertion:
\[
\begin{gathered}
 \text{For every group }G\text{ with }\operatorname{cd}_{\Z}(G)\leq2,
 \text{ and every }a=a^*\in\C[G],\\
 \left[\pi(a)\succeq0\text{ for every finite-dimensional unitary representation }\pi\text{ of }G\right]
 \Longrightarrow
 \left[a=\sum_{j=1}^{m}b_j^*b_j\text{ for some finite }m
 \text{ and }b_j\in\C[G]\right].
\end{gathered}
\]
Here \(\pi\) extends linearly to \(\C[G]\), and \(\pi(a)\succeq0\) means that the resulting Hermitian operator is positive semidefinite.

\textbf{Unresolved qualifications.} The supplied record does not establish that this displayed interpretation is the intended conjecture. It does not specify the group-ring coefficient field, the coefficients used for cohomological dimension, the representations on which positivity is tested, whether positivity is strict or non-strict, or whether the conclusion is an exact finite sum of squares. It also does not specify whether a matrix-valued version is intended. Consequently, the source question cannot be assigned a unique precise mathematical formulation from this record alone.
\subsection{Short English statement}
The question asks whether a Helton-type positivity criterion forces a sum-of-squares representation in group rings when the group has cohomological dimension at most two. The supplied record leaves the positivity criterion and several algebraic conventions unspecified, so the displayed formulation is only a possible interpretation requiring human review.
\subsection{Sources}
\begin{itemize}
\item Supplied local UnsolvedMath registry record, permanent ID \(30002899\), source problem code \texttt{catalogue entry}, in the Oberwolfach Reports Open Problems collection. The record supplies proposed year 2015 and names no individual proposer.
\item Supplied report citation: \emph{Noncommutative Geometry (2016)}, Oberwolfach Reports, pp.\ 1631--1700, DOI \texttt{10.4171/owr/2015/29}. \url{https://doi.org/10.4171/owr/2015/29}. This is the supplied authoritative source reference; its contents have not been independently inspected here.
\item Supplied literature source: ``Oberwolfach report 29,'' \url{https://doi.org/10.4171/owr/2015/29}. The literature assessment dated \texttt{2026-08-21} labels the problem open, describes the general assertion as a conjecture in the source, and reports selected known cases including \(\Z^2\) and free groups. These are supplied dated curation claims, not independently verified mathematical or literature-status conclusions.
\item The supplied background describes Oberwolfach Reports as records of mathematical workshops and their open-problem discussions, and explicitly requests independent verification of time-sensitive status and cited claims. The record also carries the supplied statement label \texttt{corrected\_verified} and a manual-cleanup verification note; these do not resolve the missing mathematical specifications identified above.
\end{itemize}
\end{document}
